\chapter[บทนำ: สถาปัตยกรรมเสมือนจริงและการคำนวณด้วย Python]{บทนำ\\สถาปัตยกรรมการจำลองเสมือนจริง AR/XR\\และการคำนวณด้วย Python}
\label{chap:introduction}

\index{WebXR}
\index{Three.js}
\index{MediaPipe}
\index{Python}
\index{การจำลองเสมือนจริง}
\index{Spatial Computing}

% ==========================================================
% แผนบริหารการสอนประจำบทที่ 1
% ==========================================================
\section*{แผนบริหารการสอนประจำบทที่ 1}
\addcontentsline{toc}{section}{แผนบริหารการสอนประจำบทที่ 1}

\begin{lessonplanbox}
\noindent\textbf{หัวข้อเนื้อหาประจำบท}
\begin{enumerate}
    \item วิวัฒนาการและบทบาทของ WebXR ในการศึกษาฟิสิกส์แผนใหม่
    \item วงจรชีวิตของ WebXR Session และการเรนเดอร์ 3 มิติ
    \item การควบคุมแบบไร้สัมผัสด้วย MediaPipe Hands
    \item วิศวกรรมลูปประสิทธิภาพสูง ไร้ขยะในหน่วยความจำ (Zero-GC Optimization)
    \item สถาปัตยกรรมการคำนวณเชิงตัวเลขและการจำลองทางวิทยาศาสตร์ด้วย Python
    \item การบูรณาการปัญญาประดิษฐ์และการเรียนรู้เชิงลึกในฟิสิกส์เชิงคำนวณ
    \item แผนผังบูรณาการเทคโนโลยีรายบท (Curriculum Technology Roadmap)
    \item สรุปเครื่องมือและเทคนิคหลักของเทคโนโลยีเสมือนจริง
\end{enumerate}

\vspace{0.4em}
\noindent\textbf{วัตถุประสงค์เชิงพฤติกรรม}
\begin{enumerate}
    \item อธิบายสถาปัตยกรรมและการทำงานร่วมกันระหว่าง WebXR Device API, Three.js และ Python ได้อย่างถูกต้อง
    \item อธิบายหลักการตรวจจับพื้นผิว (Hit-Testing) และการตรึงวัตถุ 3 มิติเชิงพื้นที่ได้อย่างชัดเจน
    \item วิเคราะห์การควบคุมท่าทางมือ 21 จุดของ MediaPipe Hands ในการมีปฏิสัมพันธ์กับแบบจำลองฟิสิกส์ได้
    \item ออกแบบการจัดการหน่วยความจำแบบ Object Pooling เพื่อรักษาอัตราเฟรมคงที่ 60 FPS ได้
    \item ประยุกต์ใช้ไลบรารี NumPy, SciPy และ SymPy ในการคำนวณเชิงตัวเลขสำหรับงานฟิสิกส์ได้อย่างเป็นระบบ
    \item อธิบายกรอบแนวคิดของ Physics-Informed Neural Networks (PINNs) และการประยุกต์ใช้ AI ในฟิสิกส์แผนใหม่ได้
\end{enumerate}

\vspace{0.4em}
\noindent\textbf{กิจกรรมการเรียนการสอน}
\begin{enumerate}
    \item การบรรยายภาพรวมเทคโนโลยีโลกเสมือนจริงและการคำนวณเชิงตัวเลขในฟิสิกส์แผนใหม่
    \item กิจกรรมสาธิตการเข้าถึงแบบจำลอง WebXR ผ่านเว็บเบราว์เซอร์บนสมาร์ตโฟนและแท็บเล็ต
    \item ปฏิบัติการทดสอบระบบตรวจจับพิกัดมือด้วย MediaPipe Hands เพื่อควบคุมตัวแปรฟิสิกส์แบบไร้สัมผัส
    \item ปฏิบัติการเขียนสคริปต์ Python พื้นฐานเพื่อเตรียมชุดข้อมูลและการแก้สมการเชิงอนุพันธ์
\end{enumerate}

\vspace{0.4em}
\noindent\textbf{สื่อการเรียนการสอน}
\begin{enumerate}
    \item เอกสารคำสอน บทที่ 1 บทนำ สถาปัตยกรรมการจำลองเสมือนจริงและการคำนวณด้วย Python
    \item สไลด์บรรยายอิเล็กทรอนิกส์และแผนภาพสถาปัตยกรรมระบบ 3 มิติ
    \item ห้องปฏิบัติการเสมือนจริงจำลองบนเว็บเบราว์เซอร์ (WebXR Demonstrator)
    \item สภาพแวดล้อมประมวลผล Python (Jupyter Notebook / Google Colab)
\end{enumerate}

\vspace{0.4em}
\noindent\textbf{การวัดและประเมินผล}
\begin{enumerate}
    \item การประเมินความสนใจและการมีส่วนร่วมในการอภิปรายและทดลองในชั้นเรียน
    \item การประเมินผลสัมฤทธิ์จากการทำใบงานและแบบฝึกหัดประจำบทที่ 1
    \item การทดสอบย่อยวัดผลสัมฤทธิ์ทางการเรียนประจำบทที่ 1
\end{enumerate}
\end{lessonplanbox}

\section{วิวัฒนาการและบทบาทของ WebXR ในการศึกษาฟิสิกส์แผนใหม่}


มโนทัศน์ในวิชาฟิสิกส์แผนใหม่ เช่น ปริภูมิ-เวลา 4 มิติ ความน่าจะเป็นของออร์บิทัลอิเล็กตรอน โครงสร้างควาร์กภายในฮาดรอน การแตกตัวของนิวเคลียส หรือความพัวพันเชิงควอนตัม เป็นปรากฏการณ์ที่ไม่สามารถมองเห็นได้ด้วยตาเปล่าในชีวิตประจำวัน และขัดแย้งกับสัญชาตญาณดั้งเดิมของมนุษย์ การจัดการเรียนรู้แบบเดิมที่พึ่งพาเพียงสมการบนกระดานดำหรือภาพวาด 2 มิติในหนังสือ จึงสร้างภาระทางพุทธิปัญญา (Cognitive Load) ให้แก่ผู้เรียนอย่างยิ่ง

การนำเทคโนโลยี \textbf{ความเป็นจริงขยาย (Extended Reality หรือ XR)} ซึ่งครอบคลุมทั้ง Augmented Reality (AR) และ Virtual Reality (VR) มาผสานเข้ากับการศึกษาวิทยาศาสตร์ จึงเปิดมิติใหม่ในการแปลงนามธรรมทางคณิตศาสตร์ให้กลายเป็นวัตถุเชิงพื้นที่ (Spatial Objects) ที่ผู้เรียนสามารถสำรวจ เดินวนรอบ สังเกตการเปลี่ยนแปลงจากทุกมุมมอง และมีปฏิสัมพันธ์แบบเรียลไทม์

\textbf{WebXR Device API} คือมาตรฐานเว็บสากลที่กำหนดขึ้นโดย World Wide Web Consortium (W3C) เพื่อเปิดทางให้เว็บเบราว์เซอร์ยุคใหม่สามารถสื่อสารโดยตรงกับฮาร์ดแวร์เซนเซอร์ตรวจจับตำแหน่งและการเคลื่อนไหว 6 แกนอิสระ (6-DoF Tracking) ของสมาร์ตโฟน แท็บเล็ต หรือแว่นตาอัจฉริยะ (Head-Mounted Displays) ข้อได้เปรียบสูงสุดของ WebXR คือการเป็นเทคโนโลยีแบบไร้รอยต่อ (Zero-Install) ผู้เรียนเพียงสแกนคิวอาร์โค้ดหรือเปิดลิงก์ผ่านเว็บเบราว์เซอร์ ก็สามารถแปลงพื้นที่จริงในห้องเรียนให้กลายเป็นห้องปฏิบัติการฟิสิกส์เสมือนจริงได้ทันที

\begin{figure}[H]
\centering
\begin{tikzpicture}[scale=0.95]
    % Real World Layer
    \draw[darkNavy, fill=gray!10, rounded corners=6pt] (-5, -1.8) rectangle (5, -0.6);
    \node at (0, -1.2) [font=\footnotesize\bfseries\color{darkNavy}] {สภาพแวดล้อมโลกจริง (Physical World \& Camera Video Feed)};
    
    % WebXR Hardware Abstraction Layer
    \draw[cyberBlue, fill=cyan!10, rounded corners=6pt] (-5, 0.0) rectangle (5, 1.2);
    \node at (0, 0.6) [font=\footnotesize\bfseries\color{cyberBlue}] {WebXR Device API (Pose, Hit-Test, View Matrix, Projection Matrix)};
    
    % Three.js & Shader Engine
    \draw[neonPurple, fill=purple!10, rounded corners=6pt] (-5, 1.8) rectangle (5, 3.0);
    \node at (0, 2.4) [font=\footnotesize\bfseries\color{neonPurple}] {Three.js Scene Graph, Custom GLSL Shaders \& WebGL/WebGPU};
    
    % Interaction Layer (MediaPipe)
    \draw[emeraldGlow, fill=green!10, rounded corners=6pt] (-5, 3.6) rectangle (5, 4.8);
    \node at (0, 4.2) [font=\footnotesize\bfseries\color{emeraldGlow!80!black}] {MediaPipe Hands 21-Joint 3D Gesture Controller (Zero-Touch Interface)};
    
    % Connecting arrows
    \draw[-{Stealth[scale=1.2]}, line width=1.5pt, slateGrey] (0, -0.6) -- (0, 0.0);
    \draw[-{Stealth[scale=1.2]}, line width=1.5pt, slateGrey] (0, 1.2) -- (0, 1.8);
    \draw[-{Stealth[scale=1.2]}, line width=1.5pt, slateGrey] (0, 3.0) -- (0, 3.6);
\end{tikzpicture}
\caption{สถาปัตยกรรมซอฟต์แวร์จำลองเสมือนจริง WebXR ร่วมกับ Three.js และ MediaPipe Hands}
\label{fig:webxr_architecture}
\end{figure}

\section{วงจรชีวิตของ WebXR Session และการเรนเดอร์ 3 มิติ}
\index{Hit-Test}
\index{WebXR Session}

หัวใจสำคัญของการทำงานร่วมกับ WebXR Device API ประกอบด้วย 4 ขั้นตอนหลัก
\begin{enumerate}
    \item \textbf{การตรวจสอบความพร้อมของอุปกรณ์ (Device Capability Check)} เบราว์เซอร์จะทำการตรวจสอบว่าระบบฮาร์ดแวร์รองรับการแสดงผลความเป็นจริงเสริมหรือไม่ ผ่านคำสั่ง \texttt{navigator.xr.isSessionSupported('immersive-ar')}
    \item \textbf{การร้องขอเปิดเซสชันเสมือนจริง (Request Session)} เมื่อผู้ใช้กดปุ่มเข้าสู่โหมดการทดลอง ระบบจะสร้างเซสชันการทำงานแบบ Immersive AR พร้อมทั้งขอสิทธิ์การเข้าถึงฟีเจอร์การตรวจจับระนาบพื้นผิว (\texttt{hit-test})
    \item \textbf{การตรวจจับพื้นผิวและการตรึงพิกัด (Hit-Testing \& Spatial Anchoring)} ในระหว่างการเคลื่อนย้ายอุปกรณ์ ระบบจะยิงลำรังสีเสมือน (Virtual Raycast) จากจุดกึ่งกลางหน้าจอไปตกกระทบระนาบพื้นผิวโลกจริง เช่น โต๊ะทดลองหรือพื้นห้อง เพื่อสร้างจุดอ้างอิงพิกัดเสมือน (Anchor Pose Matrix ขนาด $4 \times 4$) สำหรับวางอุปกรณ์ทดลองทางฟิสิกส์
    \item \textbf{วงรอบการแสดงผลแบบซิงโครนัส (Synchronous Render Loop)} ทุกครั้งที่มีการขยับกล้อง ระบบจะเรียกฟังก์ชันผ่าน \texttt{XRSession.requestAnimationFrame} เพื่อรับค่า Projection Matrix และ View Matrix ล่าสุดจากเซนเซอร์ตรวจจับการเคลื่อนไหว นำไปปรับมุมมองของกล้องเสมือนใน Three.js ให้สอดรับกับโลกจริงแบบ 1:1
\end{enumerate}

\begin{figure}[H]
\centering
\begin{tikzpicture}[scale=0.92]
    % Real table surface
    \draw[slateGrey, fill=gray!20, line width=1.2pt] (-4, -1.8) -- (4, -1.8) -- (5.5, -0.6) -- (-2.5, -0.6) -- cycle;
    \node at (1.0, -1.4) [font=\scriptsize\bfseries\color{slateGrey}] {พื้นผิวโต๊ะจริงในห้องเรียน (Physical Plane)};
    
    % Smartphone / AR Device
    \draw[darkNavy, fill=darkNavy!90, rounded corners=4pt, rotate around={-20:(-2, 2.5)}] (-3.0, 1.8) rectangle (-1.0, 3.2);
    \draw[cyan, fill=cyan!15, rounded corners=2pt, rotate around={-20:(-2, 2.5)}] (-2.9, 1.9) rectangle (-1.1, 3.1);
    \node at (-2.0, 2.8) [font=\scriptsize\bfseries\color{darkNavy}, rotate=-20] {กล้องสมาร์ตโฟน};
    
    % Virtual Raycast (Hit-test ray)
    \draw[-{Stealth[scale=1.3]}, line width=1.5pt, color=red!80!black, dashed] (-1.6, 2.1) -- (1.5, -1.2);
    \node at (0.6, 0.8) [font=\scriptsize\bfseries\color{red!80!black}, rotate=-43] {ลำแสงเสมือน (XRHitTestSource)};
    
    % Hit Point / Anchor Pose
    \filldraw[quantumGold] (1.5, -1.2) circle (4pt);
    \draw[quantumGold, ultra thick] (1.5, -1.2) circle (8pt);
    \node at (2.4, -0.8) [font=\footnotesize\bfseries\color{quantumGold!90!black}] {จุดกระทบ (Anchor Pose)};
    
    % Virtual 3D Hologram spawning on top of anchor
    \draw[cyberBlue, thick, dashed] (1.5, -1.2) -- (1.5, 0.6);
    \draw[line width=1.6pt, color=cyberBlue, fill=cyan!20] (1.5, 1.2) circle (0.6cm);
    \node at (1.5, 1.2) [font=\tiny\bfseries\color{darkNavy}] {แบบจำลอง};
    \draw[neonPurple, thick, rotate around={35:(1.5, 1.2)}] (1.5, 1.2) ellipse (1.1cm and 0.35cm);
    \node at (1.5, 2.2) [font=\footnotesize\bfseries\color{cyberBlue}] {อ็อบเจกต์ฟิสิกส์ 3 มิติ};
\end{tikzpicture}
\caption{ขั้นตอนการตรวจจับพื้นผิว (Hit-Testing) และการตรึงพิกัดเสมือนจริง (Spatial Anchoring) ใน WebXR}
\label{fig:webxr_hit_test}
\end{figure}

\section{การควบคุมแบบไร้สัมผัสด้วย MediaPipe Hands}
\index{Hand Tracking}
\index{MediaPipe}

เพื่อปลดล็อกข้อจำกัดจากการสัมผัสหน้าจอ 2 มิติ ตำราเล่มนี้ได้ผสานการทำงานของ \textbf{MediaPipe Hands} ซึ่งเป็นโมเดล Deep Learning ประสิทธิภาพสูงของ Google เข้ากับเว็บแอปพลิเคชัน โมเดลนี้มีความสามารถในการอนุมานพิกัด 3 มิติ ($x, y, z$) ของข้อต่อนิ้วมือจำนวน 21 จุด (21 3D Hand Landmarks) จากสตรีมวิดีโอของกล้องได้แบบเรียลไทม์

\begin{figure}[H]
\centering
\begin{tikzpicture}[scale=0.90]
    % Wrist
    \filldraw[darkNavy] (0, 0) circle (3pt) node[below, font=\tiny\bfseries] {0 (WRIST)};
    
    % Thumb (1, 2, 3, 4)
    \draw[thick, color=cyberBlue] (0, 0) -- (-0.8, 0.8) -- (-1.4, 1.6) -- (-1.8, 2.3) -- (-2.1, 2.9);
    \foreach \coord/\num in {(-0.8,0.8)/1, (-1.4,1.6)/2, (-1.8,2.3)/3, (-2.1,2.9)/4} {
        \filldraw[cyberBlue] \coord circle (2.5pt);
    }
    \node at (-2.3, 3.2) [font=\tiny\bfseries\color{cyberBlue}] {4 (THUMB\_TIP)};
    
    % Index finger (5, 6, 7, 8)
    \draw[thick, color=neonPurple] (0, 0) -- (-0.6, 2.0) -- (-0.8, 3.2) -- (-0.9, 4.1) -- (-1.0, 4.8);
    \foreach \coord/\num in {(-0.6,2.0)/5, (-0.8,3.2)/6, (-0.9,4.1)/7, (-1.0,4.8)/8} {
        \filldraw[neonPurple] \coord circle (2.5pt);
    }
    \node at (-1.0, 5.1) [font=\tiny\bfseries\color{neonPurple}] {8 (INDEX\_TIP)};
    
    % Middle finger (9, 10, 11, 12)
    \draw[thick, color=emeraldGlow] (0, 0) -- (0.0, 2.2) -- (0.0, 3.5) -- (0.0, 4.5) -- (0.0, 5.3);
    \foreach \coord/\num in {(0.0,2.2)/9, (0.0,3.5)/10, (0.0,4.5)/11, (0.0,5.3)/12} {
        \filldraw[emeraldGlow] \coord circle (2.5pt);
    }
    \node at (0.0, 5.6) [font=\tiny\bfseries\color{emeraldGlow}] {12 (MIDDLE\_TIP)};
    
    % Ring finger (13, 14, 15, 16)
    \draw[thick, color=quantumGold] (0, 0) -- (0.6, 2.0) -- (0.8, 3.2) -- (0.9, 4.1) -- (1.0, 4.8);
    \foreach \coord/\num in {(0.6,2.0)/13, (0.8,3.2)/14, (0.9,4.1)/15, (1.0,4.8)/16} {
        \filldraw[quantumGold] \coord circle (2.5pt);
    }
    \node at (1.0, 5.1) [font=\tiny\bfseries\color{quantumGold}] {16 (RING\_TIP)};
    
    % Pinky finger (17, 18, 19, 20)
    \draw[thick, color=red!80!black] (0, 0) -- (1.1, 1.6) -- (1.4, 2.5) -- (1.6, 3.2) -- (1.8, 3.8);
    \foreach \coord/\num in {(1.1,1.6)/17, (1.4,2.5)/18, (1.6,3.2)/19, (1.8,3.8)/20} {
        \filldraw[red!80!black] \coord circle (2.5pt);
    }
    \node at (2.0, 4.1) [font=\tiny\bfseries\color{red!80!black}] {20 (PINKY\_TIP)};
    
    % Palm connective mesh
    \draw[dashed, slateGrey] (-0.8, 0.8) -- (-0.6, 2.0) -- (0.0, 2.2) -- (0.6, 2.0) -- (1.1, 1.6);
\end{tikzpicture}
\caption{โครงข่ายข้อต่อนิ้วมือ 21 จุด 3 มิติ (MediaPipe 21-Joint 3D Hand Landmarks) สำหรับควบคุมการทดลองฟิสิกส์}
\label{fig:mediapipe_hand_landmarks}
\end{figure}

รูปแบบท่าทางมือมาตรฐาน (Standard Hand Gestures) ที่นำมาประยุกต์ใช้ในทุกบทเรียน ได้แก่
\begin{itemize}
    \item \textbf{Pinch Gesture (การบีบนิ้วโป้งและนิ้วชี้เข้าหากัน)} คำนวณจากระยะห่างระหว่างจุดที่ 4 (THUMB\_TIP) และจุดที่ 8 (INDEX\_TIP) ใช้สำหรับ "หยิบ จับ และเคลื่อนย้าย" อนุภาค อะตอม หรือกรอบอ้างอิงในมิติ 3 มิติ
    \item \textbf{Palm Facing (การหงายฝ่ามือขึ้น)} ใช้เวกเตอร์แนวฉากของฝ่ามือ (Normal Vector) เพื่อสร้างแท่นรองรับอนุภาคเสมือนจริง เช่น ใช้เป็นแท่นตรวจจับรังสี หรือแท่นชนอนุภาค
    \item \textbf{Index Pointing (การชี้นิ้วชี้)} ใช้เวกเตอร์ทิศทางจากจุดที่ 5 ไปยังจุดที่ 8 เพื่อใช้เป็นลำเลเซอร์เสมือนในการเลือกสถานะพลังงานควอนตัม ($n$) หรือปรับสเกลค่าพารามิเตอร์ทางกายภาพ
\end{itemize}

\section{วิศวกรรมลูปประสิทธิภาพสูง ไร้ขยะในหน่วยความจำ}
\index{Zero-GC}
\index{Garbage Collection}

ในการเรนเดอร์กราฟิกเสมือนจริงแบบเรียลไทม์ที่ความถี่ 60 เฟรมต่อวินาที (60 FPS) เบราว์เซอร์มีช่วงเวลาจำกัดเพียง $16.67\text{ มิลลิวินาที}$ ต่อหนึ่งเฟรม ในการประมวลผลฟิสิกส์ การคำนวณตำแหน่ง และการเรนเดอร์ภาพ หากนักพัฒนาเขียนโค้ดที่สร้างอ็อบเจกต์ใหม่ในทุกเฟรม เช่น คำสั่ง \texttt{new THREE.Vector3()} ขยะในหน่วยความจำ (Garbage) จะถูกสะสมอย่างรวดเร็ว ส่งผลให้เบราว์เซอร์ต้องหยุดการทำงานชั่วขณะเพื่อทำความสะอาดหน่วยความจำ (Garbage Collection Pause) ทำให้ภาพจำลองกระตุกและเกิดอาการเมารถเสมือนจริง (Simulator Sickness)

\begin{pitfallbox}[กฎเหล็กวิศวกรรมลูปแอนิเมชัน 60 FPS (Zero-GC Loop Engineering)]
\textbf{แนวทางปฏิบัติเพื่อประสิทธิภาพสูงสุดในการพัฒนาสื่อฟิสิกส์เสมือนจริง}
\begin{itemize}
    \item \textbf{หลีกเลี่ยงการสร้างอินสแตนซ์ใหม่ในลูป:} ห้ามใช้คำสั่ง \texttt{new} ภายในเมธอด \texttt{renderLoop} หรือ \texttt{requestAnimationFrame} โดยเด็ดขาด
    \item \textbf{ใช้หน่วยความจำชั่วคราวส่วนกลาง (Pre-allocated Scratchpad):} ประกาศเวกเตอร์และเมทริกซ์ชั่วคราวไว้นอกลูปหลัก แล้วใช้เมธอด \texttt{.copy()}, \texttt{.set()}, หรือ \texttt{.add()} ในการคำนวณซ้ำ
    \item \textbf{การคืนทรัพยากรหน่วยความจำกราฟิก:} เมื่อเสร็จสิ้นการทดลอง ต้องเรียกใช้ฟังก์ชัน \texttt{.dispose()} บน Geometry และ Material ทุกชิ้น เพื่อป้องกันหน่วยความจำ VRAM รั่วไหล (VRAM Leaks)
\end{itemize}
\end{pitfallbox}

\section{สถาปัตยกรรมการคำนวณเชิงตัวเลขและการจำลองทางวิทยาศาสตร์ด้วย Python}
\index{NumPy}
\index{SciPy}
\index{Matplotlib}

ควบคู่ไปกับการเรนเดอร์ 3 มิติเชิงพื้นที่ การคำนวณเชิงตัวเลข (Scientific Numerical Computing) คือเสาหลักที่สองของตำราเล่มนี้ ภาษา Python ได้รับการคัดเลือกให้เป็นเครื่องมือหลักในการแก้ปัญหาทางฟิสิกส์ เนื่องจากมีระบบนิเวศของไลบรารีทางวิทยาศาสตร์ที่สมบูรณ์และได้มาตรฐานระดับสากล

เครื่องมือหลักของ Python ที่ใช้ในตำราเล่มนี้ประกอบด้วย
\begin{enumerate}
    \item \textbf{NumPy (Numerical Python)} โครงสร้างข้อมูลแบบอาเรย์หลายมิติ (N-dimensional Array) ที่ประมวลผลบนหน่วยความจำแบบต่อเนื่องด้วยภาษา C ช่วยให้การคำนวณเวกเตอร์และเมทริกซ์ในทฤษฎีสัมพัทธภาพและกลศาสตร์ควอนตัมทำได้อย่างรวดเร็ว
    \item \textbf{SciPy (Scientific Python)} โมดูลการคำนวณทางวิทยาศาสตร์ชั้นสูง โดยเฉพาะ \texttt{scipy.integrate.solve\_ivp} สำหรับการแก้ระบบสมการเชิงอนุพันธ์สามัญ (ODEs) ในการจำลองการสลายตัวของนิวเคลียสและแบบจำลองจักรวาลวิทยา ตลอดจน \texttt{scipy.linalg} สำหรับการหาค่าเจาะจง (Eigenvalues) ของฮามิลโทเนียนควอนตัม
    \item \textbf{Matplotlib} การพล็อตและแสดงผลข้อมูลเชิงวิทยาศาสตร์ ทั้งกราฟ 2 มิติระดับตีพิมพ์ และแผนผังการกระจายตัวของความน่าจะเป็นเชิงควอนตัม
    \item \textbf{SymPy (Symbolic Python)} การแก้สมการเชิงสัญลักษณ์ การแปลงพีชคณิต และการหาอนุพันธ์และปริพันธ์เชิงวิเคราะห์เพื่อตรวจสอบความถูกต้องของทฤษฎี
\end{enumerate}

\section{การบูรณาการปัญญาประดิษฐ์และการเรียนรู้เชิงลึกในฟิสิกส์เชิงคำนวณ}
\index{ปัญญาประดิษฐ์ในฟิสิกส์}
\index{การเรียนรู้เชิงลึก}
\index{PINN}
\index{Physics-Informed Neural Networks}

พัฒนาการล่าสุดในวงการวิทยาศาสตร์คำนวณยุคปัจจุบันคือการผสาน \textbf{ปัญญาประดิษฐ์ (Artificial Intelligence: AI)} และ \textbf{การเรียนรู้เชิงลึก (Deep Learning)} เข้ากับการจำลองทางฟิสิกส์ ซึ่งนำไปสู่วิทยาการแขนงใหม่ที่เรียกว่า \textit{การเรียนรู้ของเครื่องเชิงวิทยาศาสตร์ (Scientific Machine Learning: SciML)} หรือ \textit{AI for Science}

กระบวนทัศน์ดั้งเดิมของการเรียนรู้เชิงลึกมักถูกวิจารณ์ว่าเป็น "กล่องดำ (Black Box)" ที่เรียนรู้ความสัมพันธ์ของข้อมูลโดยไม่เข้าใจกฎเกณฑ์ธรรมชาติ แต่ในงานฟิสิกส์แผนใหม่ ได้มีการคิดค้นสถาปัตยกรรม \textbf{โครงข่ายประสาทเทียมที่รู้ฟิสิกส์ (Physics-Informed Neural Networks: PINNs)} ซึ่งพัฒนาขึ้นโดย Raissi และคณะ (2019) โดยมีหัวใจสำคัญคือการฝังกฎทางฟิสิกส์ เช่น สมการเชิงอนุพันธ์ย่อย (PDEs) และกฎการอนุรักษ์ เข้าไปในฟังก์ชันการสูญเสีย (Loss Function) ของโครงข่ายประสาทเทียมโดยตรง

\begin{equation}
\mathcal{L}_{\text{total}} = \mathcal{L}_{\text{data}} + \lambda_{\text{phys}} \mathcal{L}_{\text{physics}}
\end{equation}

เมื่อ $\mathcal{L}_{\text{data}}$ คือผลต่างระหว่างการทำนายของโครงข่ายประสาทกับข้อมูลที่สังเกตได้ $\mathcal{L}_{\text{physics}}$ คือส่วนตกค้าง (Residual) ของสมการฟิสิกส์ เช่น สมการชเรอดิงเงอร์ หรือสมการการเคลื่อนที่สัมพัทธภาพ และ $\lambda_{\text{phys}}$ คือค่าน้ำหนักควบคุมดุลยภาพ วิธีการนี้ช่วยให้โมเดลสามารถพยากรณ์พฤติกรรมของระบบควอนตัมและสัมพัทธภาพได้อย่างแม่นยำแม้มีข้อมูลตรวจวัดจำกัด

นอกจากนี้ การผนวกรวมโมเดลการเรียนรู้เชิงลึกแบบย่อส่วน (Surrogate Models) เข้ากับระบบแสดงผล WebXR บนเว็บเบราว์เซอร์ ยังช่วยลดภาระการคำนวณที่ซับซ้อน ทำให้สามารถเรนเดอร์สนามควอนตัมและวิถีการเคลื่อนที่สัมพัทธภาพได้แบบเรียลไทม์ที่ 60 เฟรมต่อวินาที ซึ่งจะมีการนำเสนอตัวอย่างและการประยุกต์ใช้ในแต่ละบทเรียนต่อไป

\section{แผนผังบูรณาการเทคโนโลยีรายบท}

เพื่อให้นักศึกษาและอาจารย์ผู้สอนเห็นภาพรวมอย่างชัดเจนว่าจะมีการนำเทคโนโลยี AR/XR และ Python ไปประยุกต์ใช้อย่างไรในเนื้อหาแต่ละบท ตารางที่ \ref{tab:curriculum_tech_map} แสดงแผนผังการกระจายเทคโนโลยีตลอดทั้งเล่ม

\begin{table}[H]
\centering
\caption{แผนผังการบูรณาการเทคโนโลยีเสมือนจริง AR/XR และการคำนวณ Python ในตำราฟิสิกส์แผนใหม่}
\label{tab:curriculum_tech_map}
\small
\begin{tabularx}{\textwidth}{c p{3.2cm} Y Y}
\toprule
\textbf{บทที่} & \textbf{ชื่อบทเรียน} & \textbf{การจำลองเสมือนจริง AR/XR (WebXR)} & \textbf{การคำนวณเชิงตัวเลขด้วย Python} \\
\midrule
2 & ทฤษฎีสัมพัทธภาพพิเศษและเรขาคณิตกาล-อวกาศ & การจำลอง 4D Spacetime กรวยแสง และการปรับพิกัด LERP Smoothing & การวิเคราะห์ตัวคูณลอเรนซ์ โมเมนตัม และพลังงานเชิงสัมพัทธภาพ \\
\midrule
3 & กลศาสตร์ควอนตัมและโครงสร้างอะตอม & แบบจำลองอะตอมบอร์เสมือนจริง WebXR โค้ดเต็ม และออร์บิทัล 3 มิติ & ระเบียบวิธีผลต่างอันตะ (Finite Difference) แก้สมการชเรอดิงเงอร์ \\
\midrule
4 & ฟิสิกส์ของอนุภาคและแบบจำลองมาตรฐาน & โครงสร้างควาร์กในแบริออน-มีซอน และการตรึงพิกัดด้วย Hit-Test & การประกอบมวลไม่แปรเปลี่ยน (Invariant Mass) ของอนุภาคฮิกส์ \\
\midrule
5 & ฟิสิกส์นิวเคลียร์และรังสี & การจำลองปฏิกิริยาลูกโซ่นิวเคลียร์ฟิชชันแบบ Cascade 3 มิติ & การแก้ระบบสมการลูกโซ่การสลายตัวของเบตแมนด้วย \texttt{solve\_ivp} \\
\midrule
6 & สารสนเทศและการคำนวณเชิงควอนตัม & ทรงกลมบล็อค 3 มิติแบบมีปฏิสัมพันธ์ และการหมุนเวกเตอร์สถานะ & ตัวจำลองเวกเตอร์สถานะควอนตัม (Quantum Statevector Simulator) \\
\midrule
7 & ดาราศาสตร์ฟิสิกส์และจักรวาลวิทยา & กาล-อวกาศโค้งรอบหลุมดำ เลนส์ความโน้มถ่วง และขอบฟ้าเหตุการณ์ & การแก้สมการฟรีดมันน์เพื่อจำลองการขยายตัวของเอกภพตามยุคสมัย \\
\bottomrule
\end{tabularx}
\end{table}

\section{สรุปเครื่องมือและเทคนิคหลักของเทคโนโลยีเสมือนจริง}

\begin{table}[H]
\centering
\caption{สรุปเครื่องมือและเทคนิคหลักในการสร้างแบบจำลองฟิสิกส์ด้วย WebXR}
\label{tab:ar_xr_summary_intro}
\small
\begin{tabularx}{\textwidth}{l >{\centering\arraybackslash}p{3.8cm} Y}
\toprule
\textbf{เทคโนโลยี / ฟังก์ชัน} & \textbf{เครื่องมือหลัก} & \textbf{หน้าที่ในการศึกษาฟิสิกส์} \\
\midrule
WebXR Device API & W3C Standard & เชื่อมต่อฮาร์ดแวร์ AR/VR บนเว็บเบราว์เซอร์โดยตรงโดยไม่ต้องติดตั้งแอปพลิเคชัน \\
Three.js & WebGL / WebGPU & จัดการฉาก เรนเดอร์วัสดุ แสงเงาแบบ PBR และตาข่ายสามมิติ \\
MediaPipe Hands & Machine Learning & ติดตามข้อต่อนิ้วมือ 21 จุด เพื่อควบคุมการทดลองแบบไร้การสัมผัส \\
Zero-GC Loop & Object Pooling & ป้องกันการกระตุกของภาพ รักษาอัตราเฟรมคงที่ 60 FPS \\
Spatial LERP & $(1-\alpha)\mathbf{x}_0 + \alpha\mathbf{x}_1$ & ลดสัญญาณรบกวนจากการตรวจจับตำแหน่งของเซนเซอร์กล้อง AR \\
Hit-Testing & Raycasting & ตรวจจับระนาบโต๊ะทดลองในโลกจริงเพื่อตรึงแบบจำลอง 3 มิติ \\
\bottomrule
\end{tabularx}
\end{table}

\newpage
